% p101_1; source PDF page 101; manually verified editable table.
\begin{tikzpicture}[x=1cm,y=-1cm,inner sep=0pt,outer sep=0pt]
\path[use as bounding box] (0,0) rectangle (16.000,7.865);
\draw[line width=.4pt] (0.000,0.000) rectangle (4.000,0.855);
\node[align=center,text width=3.740cm,font={\linespread{1}\fontsize{10}{13.50}\selectfont\bfseries}] at (2.000,0.427) {偶然避险的辨认标准};
\draw[line width=.4pt] (4.000,0.000) rectangle (7.200,0.855);
\node[align=center,text width=2.940cm,font={\linespread{1}\fontsize{10}{13.50}\selectfont\bfseries}] at (5.600,0.427) {范例};
\draw[line width=.4pt] (7.200,0.000) rectangle (11.800,0.855);
\node[align=center,text width=4.340cm,font={\linespread{1}\fontsize{10}{13.50}\selectfont\bfseries}] at (9.500,0.427) {观点展示};
\draw[line width=.4pt] (11.800,0.000) rectangle (16.000,0.855);
\node[align=center,text width=3.940cm,font={\linespread{1}\fontsize{10}{13.50}\selectfont\bfseries}] at (13.900,0.427) {处理结论};
\draw[line width=.4pt] (0.000,0.855) rectangle (4.000,7.865);
\node[align=left,text width=3.740cm,font={\linespread{1}\fontsize{10}{13.50}\selectfont}] at (2.000,4.360) {在避免危险这一点上，没有做到主客观相一致，也即客观上避免了危险，但行为人主观上没有认识到这一点};
\draw[line width=.4pt] (4.000,0.855) rectangle (7.200,7.865);
\node[align=left,text width=2.940cm,font={\linespread{1}\fontsize{10}{13.50}\selectfont}] at (5.600,4.360) {甲以毁坏财物的目的砸坏乙的汽车玻璃，却碰巧救了被锁在车里的快要窒息的小孩};
\draw[line width=.4pt] (7.200,0.855) rectangle (11.800,4.360);
\node[align=left,text width=4.340cm,font={\linespread{1}\fontsize{10}{13.50}\selectfont}] at (9.500,2.607) {避险认识不要说：成立紧急避险，只要求客观上避免了危险，不要求主观上认识到这一点。结论：偶然避险成立紧急避险};
\draw[line width=.4pt] (11.800,0.855) rectangle (16.000,4.360);
\node[align=left,text width=3.940cm,font={\linespread{1}\fontsize{10}{13.50}\selectfont}] at (13.900,2.607) {甲成立紧急避险，不构成故意毁坏财物罪};
\draw[line width=.4pt] (7.200,4.360) rectangle (11.800,7.865);
\node[align=left,text width=4.340cm,font={\linespread{1}\fontsize{10}{13.50}\selectfont}] at (9.500,6.112) {避险认识必要说：成立紧急避险，要求主观上具有避险认识。结论：偶然避险不成立紧急避险};
\draw[line width=.4pt] (11.800,4.360) rectangle (16.000,7.865);
\node[align=left,text width=3.940cm,font={\linespread{1}\fontsize{10}{13.50}\selectfont}] at (13.900,6.112) {甲不成立紧急避险，成立故意毁坏财物罪，但是由于甲制造了好结果，所以不构成故意毁坏财物罪既遂，而构成未遂};
\end{tikzpicture}
